\documentclass[11pt,a4paper]{article}

% ===================== ENCODAGE & LANGUE =====================
\usepackage[utf8]{inputenc}
\usepackage[T1]{fontenc}
\usepackage[french]{babel}
\usepackage{lmodern}
\usepackage{textcomp}
\DeclareUnicodeCharacter{20AC}{\texteuro}
\newcommand{\eur}{\texteuro}

% ===================== MISE EN PAGE =====================
\usepackage[a4paper,margin=2.3cm]{geometry}
\usepackage{xcolor}
\usepackage{graphicx}
\usepackage{booktabs}
\usepackage{array}
\usepackage{enumitem}
\usepackage{listings}
\usepackage{titlesec}
\usepackage{fancyhdr}
\usepackage{tikz}
\usetikzlibrary{shapes.geometric, arrows.meta, positioning, fit, backgrounds, calc}
\usepackage[hidelinks]{hyperref}

\setlength{\parindent}{0pt}
\setlength{\parskip}{0.45em}

% ===================== COULEURS DE LA CHARTE GEFRADIS =====================
\definecolor{gefblue}{HTML}{253081}     % St Patricks blue
\definecolor{gefmidnight}{HTML}{040C4D}  % Midnight blue
\definecolor{gefyellow}{HTML}{FFD52B}    % Cyber yellow
\definecolor{gefbeige}{HTML}{FFFDE5}     % Beige
\definecolor{gefghost}{HTML}{E7E9F4}     % Ghost white
\definecolor{gefgrey}{HTML}{E4E4E4}      % Platinum

% ===================== STYLES DE TITRES =====================
\titleformat{\section}{\Large\bfseries\sffamily\color{gefblue}}{\thesection.}{0.6em}{}
\titleformat{\subsection}{\large\bfseries\sffamily\color{gefmidnight}}{\thesubsection.}{0.6em}{}
\titlespacing*{\section}{0pt}{1.4em}{0.6em}

% ===================== EN-TETE / PIED =====================
\setlength{\headheight}{22pt}
\addtolength{\topmargin}{-10pt}
\pagestyle{fancy}
\fancyhf{}
\fancyhead[L]{\small\sffamily\color{gefblue}Cahier de conception}
\fancyhead[R]{\small\sffamily Générateur de kits graphiques}
\fancyfoot[C]{\small\sffamily\thepage}
\renewcommand{\headrulewidth}{0.6pt}
\renewcommand{\headrule}{\hbox to\headwidth{\color{gefblue}\leaders\hrule height \headrulewidth\hfill}}

% ===================== CODE / JSON =====================
\lstset{
  basicstyle=\ttfamily\footnotesize,
  frame=single,
  rulecolor=\color{gefblue},
  backgroundcolor=\color{gefghost!45},
  breaklines=true,
  columns=fullflexible,
  showstringspaces=false,
  xleftmargin=4pt,
  framexleftmargin=2pt,
}

% ===================== STYLES TIKZ =====================
\tikzset{
  >={Stealth[round]},
  block/.style={rectangle, rounded corners=2pt, draw=gefblue, very thick, fill=gefblue!8,
                align=center, minimum width=2.3cm, minimum height=1cm, font=\small\sffamily},
  store/.style={rectangle, draw=gefmidnight, very thick, fill=gefyellow!30,
                align=center, minimum width=2.3cm, minimum height=1cm, font=\small\sffamily},
  ext/.style={rectangle, rounded corners=2pt, draw=black!55, very thick, fill=black!4,
              align=center, minimum width=2.3cm, minimum height=1cm, font=\small\sffamily},
  uc/.style={ellipse, draw=gefblue, very thick, fill=gefblue!7,
             align=center, minimum width=3.6cm, minimum height=1.05cm, font=\footnotesize\sffamily, inner sep=1pt},
  proc/.style={rectangle, rounded corners=2pt, draw=gefblue, very thick, fill=gefblue!8,
               align=center, text width=2.6cm, minimum height=0.9cm, font=\footnotesize\sffamily},
  decision/.style={diamond, draw=gefblue, very thick, fill=gefyellow!30,
                   align=center, aspect=2.9, inner sep=0pt, font=\scriptsize\sffamily},
  term/.style={rectangle, rounded corners=9pt, draw=gefmidnight, very thick, fill=gefghost,
               align=center, minimum width=2.2cm, minimum height=0.8cm, font=\footnotesize\sffamily},
  arr/.style={->, very thick, gefmidnight},
  flow/.style={->, very thick, gefblue},
}

% ===================== PAGE DE TITRE =====================
\begin{document}
\begin{titlepage}
\thispagestyle{empty}
\vspace*{1.5cm}
{\color{gefblue}\rule{\textwidth}{3pt}}\\[0.3cm]
{\color{gefyellow}\rule{\textwidth}{8pt}}\\[1.5cm]

\begin{flushleft}
{\sffamily\Huge\bfseries\color{gefblue} Cahier de conception}\\[0.6cm]
{\sffamily\LARGE\color{gefmidnight} Outil de génération automatique\\ de kits graphiques d'opérations commerciales}\\[1.0cm]
{\sffamily\large Test de développement : alternance Développeur IA}\\[0.2cm]
{\sffamily\large Société \textbf{Gefradis} (groupe Im4Pulse)}\\[2.0cm]
\end{flushleft}

{\color{gefyellow}\rule{\textwidth}{8pt}}\\[0.3cm]
{\color{gefblue}\rule{\textwidth}{3pt}}\\[1.2cm]

\begin{flushleft}
{\sffamily\large
\textbf{Auteur :} Vaneck \\[0.2cm]
\textbf{Version :} 0.3 (conception)\\[0.2cm]
\textbf{Date :} \today
}
\end{flushleft}

\vfill
{\small\itshape Document technique décrivant la conception de l'application
avant le développement. Le cahier des charges (le besoin) est défini par la note
de cadrage fournie par le client ; le présent document décrit la solution retenue.}
\end{titlepage}

% ===================== SOMMAIRE =====================
\tableofcontents
\newpage

% =====================================================================
\section{Contexte et objectif}
% =====================================================================

\subsection{Contexte}
Gefradis est une entreprise de menuiserie PVC sur-mesure (portes, fenêtres),
qui vend en ligne aux particuliers et aux professionnels. Le service marketing
mène régulièrement des \textbf{opérations commerciales}. Aujourd'hui,
la production des visuels (les bannières) de ces opérations dépend de ressources
créatives et graphiques, ce qui ralentit le déploiement des campagnes.

\subsection{Objectif}
Mettre à disposition du métier marketing un \textbf{outil web} qui lui permet de
réaliser, de manière autonome et sans graphiste, un \textbf{kit graphique}
complet (un ensemble de bannières à tous les formats utiles), conforme à la
charte Gefradis, prêt à être déployé sur le site et sur les plateformes
publicitaires (Google et Meta).

\subsection{Périmètre}
\begin{itemize}[leftmargin=1.4em]
  \item \textbf{Dans le périmètre :} la saisie d'un brief, la génération du texte
        des bannières par IA, la composition des bannières à la charte (tous les
        formats de la note), l'aperçu et l'ajustement avant l'export, le
        téléchargement du kit. Versions statique et animée.
  \item \textbf{Hors périmètre (pour ce rendu) :} la gestion des comptes
        utilisateurs, la connexion directe aux régies publicitaires (publication
        automatique) et un éditeur graphique manuel libre.
\end{itemize}

% =====================================================================
\section{Besoins fonctionnels}
% =====================================================================
L'outil suit le principe \textbf{entrée, traitement, sortie}.

\begin{itemize}[leftmargin=1.4em]
  \item \textbf{F1. Saisie du brief :} l'utilisateur décrit la mécanique
        commerciale (1.1), le message commercial (1.2) et, en option, les
        illustrations (1.3 : style souhaité et demande éventuelle d'une image).
  \item \textbf{F2. Sélection des formats :} l'utilisateur coche les bannières
        à produire (site, Google Display, Google Pmax, Meta).
  \item \textbf{F3. Génération du texte :} le système transforme le brief en
        contenu structuré (accroche, offre de type montant, pourcentage ou texte,
        appel à l'action) et déduit la catégorie, la couleur de fond et le besoin
        d'image.
  \item \textbf{F4. Composition :} selon le \textbf{type} de bannière (bandeau site,
        header, publicité, format fin), le système assemble le texte, les couleurs,
        la police, le logo et, si besoin, l'image de la catégorie (banque d'images),
        aux dimensions exactes et à la charte.
  \item \textbf{F5. Aperçu :} le système affiche toutes les bannières générées.
  \item \textbf{F6. Ajustement et régénération :} l'utilisateur peut modifier
        le brief et relancer la génération autant de fois que nécessaire.
  \item \textbf{F7. Téléchargement :} le système livre le kit complet sous la
        forme d'une archive \texttt{.zip}.
\end{itemize}

% =====================================================================
\section{Besoins non fonctionnels}
% =====================================================================
\begin{itemize}[leftmargin=1.4em]
  \item \textbf{NF1. Conformité à la charte :} respect strict et systématique
        des couleurs, de la police (Cabin), du logo et du ton de la marque.
  \item \textbf{NF2. Simplicité :} l'outil doit être utilisable par un profil non
        technique, sans logiciel de graphisme.
  \item \textbf{NF3. Rapidité :} la génération d'un kit s'effectue en un temps
        raisonnable (de quelques secondes à quelques dizaines de secondes).
  \item \textbf{NF4. Maîtrise des coûts :} usage raisonné de l'API (le modèle ne
        génère que du texte court).
  \item \textbf{NF5. Sécurité :} la clé d'API n'est jamais écrite dans le code
        (fichier \texttt{.env} exclu du dépôt).
  \item \textbf{NF6. Reproductibilité :} environnement isolé (venv) décrit par un
        fichier de dépendances.
\end{itemize}

% =====================================================================
\section{Catalogue des formats à produire}
% =====================================================================
Les dimensions du site sont à confirmer auprès du client ; celles des régies
publicitaires sont des standards.

\begin{center}\small
\begin{tabular}{@{}>{\raggedright\arraybackslash}p{3.0cm} >{\raggedright\arraybackslash}p{4.3cm} >{\raggedright\arraybackslash}p{4.5cm}@{}}
\toprule
\textbf{Destination} & \textbf{Élément} & \textbf{Dimensions (px) / type} \\
\midrule
Site & Header page d'accueil & Desktop et mobile (à préciser) \\
Site & Bandeau page catégorie & Desktop et mobile (à préciser) \\
Google Display & Pavé / Bannière / Skyscraper / Mobile & 300x250, 728x90, 160x600, 320x50 (statique et animé) \\
Google Pmax & Visuels statiques & Paysage 1200x628, Carré 1200x1200, Portrait 960x1200 \\
Meta (FB / Insta) & Carré / Portrait / Story & 1080x1080, 1080x1350, 1080x1920 \\
\bottomrule
\end{tabular}
\end{center}

% =====================================================================
\section{Spécification graphique (charte)}
% =====================================================================
Ces éléments constituent les \emph{jetons de design} (design tokens) appliqués
par le moteur de composition.

\subsection{Couleurs}
\begin{center}\small
\begin{tabular}{@{}l l l@{}}
\toprule
\textbf{Nom} & \textbf{Hex} & \textbf{Usage} \\
\midrule
St Patricks blue & \texttt{\#253081} & Fond principal des bannières \\
Midnight blue & \texttt{\#040C4D} & Variante foncée \\
Cyber yellow & \texttt{\#FFD52B} & Accent et mise en avant du chiffre clé \\
Beige & \texttt{\#FFFDE5} & Fond clair / arrière-plan \\
Blanc & \texttt{\#FFFFFF} & Texte sur fond bleu \\
\bottomrule
\end{tabular}
\end{center}

\subsection{Typographie, logo, ton}
\begin{itemize}[leftmargin=1.4em]
  \item \textbf{Police :} Cabin (Regular et Bold). Elle est \emph{self-hosted}
        (fichier \texttt{.ttf} embarqué dans le projet via \texttt{@font-face}),
        avec une police sans-serif système en repli. Pas de dépendance à un CDN au
        moment du rendu.
  \item \textbf{Logo :} deux versions (bleu sur fond clair ; blanc sur fond
        foncé), à poser selon la couleur du fond.
  \item \textbf{Ton :} professeur ou mentor, précis, professionnel, clair sans
        être trop technique.
  \item \textbf{Imagerie :} esthétique épurée, minimaliste, « sortie d'usine ».
\end{itemize}


% =====================================================================
\section{Architecture}
% =====================================================================

\subsection{Architecture fonctionnelle (exécution locale)}
La figure~\ref{fig:archi} montre les composants et le flux principal.

\begin{figure}[h]
\centering
\resizebox{\textwidth}{!}{%
\begin{tikzpicture}[node distance=0.95cm and 1.25cm]
  \node[ext] (user) {Responsable\\Marketing};
  \node[block, right=of user] (front) {Frontend\\HTML/Tailwind};
  \node[block, right=of front] (back) {Backend\\FastAPI};
  \node[block, right=of back] (compo) {Moteur de\\composition};
  \node[store, right=of compo] (kit) {Kit\\(.zip)};
  \node[block, above=0.9cm of back] (copy) {Module Copy\\(Claude)};
  \node[store, below=0.9cm of compo] (assets) {Assets de marque\\(couleurs, Cabin, logo)};

  \draw[arr] (user) -- (front);
  \draw[arr] (front) -- node[above,font=\scriptsize]{HTTP} (back);
  \draw[arr, <->] (back) -- node[right,font=\scriptsize]{brief / JSON} (copy);
  \draw[arr] (back) -- node[above,font=\scriptsize]{texte} (compo);
  \draw[arr] (compo) -- node[above,font=\scriptsize]{PNG} (kit);
  \draw[arr] (assets) -- (compo);
  \draw[arr, dashed] (kit.south) to[bend right=22] node[below,font=\scriptsize]{aperçu} (front.south);
\end{tikzpicture}%
}
\caption{Architecture fonctionnelle (exécution locale).}
\label{fig:archi}
\end{figure}

\subsection{Architecture de production (GCP)}
En production, l'application est hébergée et appelle les modèles via les services
Google Cloud (figure~\ref{fig:gcp}).

\begin{figure}[h]
\centering
\begin{tikzpicture}[node distance=1.0cm and 1.5cm, font=\small\sffamily]
  \node[ext] (user) {Marketing};
  \node[block, right=2.0cm of user, minimum width=3.4cm] (app) {App FastAPI\\(+ frontend)};
  \node[draw=gefblue, very thick, dashed, rounded corners, fit=(app), inner sep=10pt, label={[gefblue,font=\footnotesize\bfseries]above:Cloud Run}] (cr) {};
  \node[block, right=3.0cm of app] (vertex) {Vertex AI\\(Claude)};
  \node[store, below=1.2cm of app] (storage) {Cloud Storage\\(assets + kits)};

  \draw[arr] (user) -- node[above,font=\scriptsize]{URL} (app);
  \draw[arr] (app) -- node[above,font=\scriptsize]{copy} (vertex);
  \draw[arr] (app) -- (storage);
\end{tikzpicture}
\caption{Architecture cible en production sur Google Cloud Platform.}
\label{fig:gcp}
\end{figure}

\subsection{Pile technique}
\begin{center}\small
\begin{tabular}{@{}l l@{}}
\toprule
\textbf{Couche} & \textbf{Technologie} \\
\midrule
Langage & Python \\
Backend / API & FastAPI \\
Frontend (interface) & HTML/CSS et Tailwind \\
Génération de texte & API Claude (SDK \texttt{anthropic}) \\
Composition d'images & HTML/CSS rendu en PNG via Playwright \\
Animations & GIF, HTML5, MP4 (à partir d'une source unique) \\
Production & Cloud Run, Vertex AI, Cloud Storage \\
\bottomrule
\end{tabular}
\end{center}

% =====================================================================
\section{Cas d'utilisation}
% =====================================================================
La figure~\ref{fig:uc} présente le diagramme de cas d'utilisation. L'acteur
unique est le \textbf{Responsable Marketing}.

\begin{figure}[h]
\centering
\begin{tikzpicture}[font=\footnotesize\sffamily]
  % --- acteur (bonhomme bâton) ---
  \draw[very thick] (0,0.55) circle (0.2);
  \draw[very thick] (0,0.35) -- (0,-0.25);
  \draw[very thick] (-0.32,0.1) -- (0.32,0.1);
  \draw[very thick] (0,-0.25) -- (-0.26,-0.7);
  \draw[very thick] (0,-0.25) -- (0.26,-0.7);
  \node[align=center, font=\footnotesize\bfseries] at (0,-1.15) {Responsable\\Marketing};

  % --- frontière du système ---
  \draw[gefblue, very thick, rounded corners]
        (2.6,-4.7) rectangle (9.6,3.3);
  \node[gefblue, font=\small\bfseries, anchor=north] at (6.1,3.25)
        {Outil de génération de kits graphiques};

  % --- cas d'utilisation ---
  \node[uc] (u1) at (6.1,2.2)  {Saisir le brief};
  \node[uc] (u2) at (6.1,1.05) {Sélectionner les formats};
  \node[uc] (u3) at (6.1,-0.1) {Générer le kit};
  \node[uc] (u4) at (6.1,-1.25){Prévisualiser les bannières};
  \node[uc] (u5) at (6.1,-2.4) {Ajuster et régénérer};
  \node[uc] (u6) at (6.1,-3.55){Télécharger le kit};

  % --- associations ---
  \foreach \u in {u1,u2,u3,u4,u5,u6}{
    \draw[very thick, gefmidnight] (0.45,0.05) -- (\u.west);
  }
\end{tikzpicture}
\caption{Diagramme de cas d'utilisation.}
\label{fig:uc}
\end{figure}

\textbf{Scénario nominal} (\emph{Générer le kit}) : l'utilisateur saisit le brief
et sélectionne les formats, puis déclenche la génération ; le système produit les
bannières et les affiche ; l'utilisateur prévisualise, ajuste si besoin (la
boucle), puis télécharge le kit.

% =====================================================================
\section{Vue dynamique : séquence de génération}
% =====================================================================
La figure~\ref{fig:seq} détaille les échanges lors d'une génération de kit, y
compris la \textbf{visualisation de l'aperçu} et la \textbf{boucle de
régénération} : l'utilisateur ajuste son brief et relance autant de fois que
nécessaire (cadre \emph{loop}) avant de télécharger le kit.

\begin{figure}[h]
\centering
\resizebox{\textwidth}{!}{%
\begin{tikzpicture}[font=\footnotesize\sffamily]
  % objets
  \node[block] (U)  at (0,0)    {Utilisateur};
  \node[block] (UI) at (4.0,0)  {Application\\(FastAPI)};
  \node[block] (CL) at (8.0,0)  {API Claude};
  \node[block] (MO) at (12.0,0) {Moteur\\composition};
  % lignes de vie
  \foreach \n in {U,UI,CL,MO}{ \draw[dashed, gray] (\n.south) -- ++(0,-10.6); }
  % message initial
  \draw[arr] (0,-1.0)  -- (4.0,-1.0)  node[midway,above,font=\scriptsize]{saisit le brief};
  % --- cadre de boucle (loop UML) ---
  \draw[gefblue, thick] (-1.1,-1.7) rectangle (13.3,-8.3);
  \fill[gefblue] (-1.1,-1.7) rectangle (-0.2,-2.2);
  \node[white, font=\scriptsize\bfseries] at (-0.65,-1.95) {loop};
  \node[gefmidnight, font=\scriptsize, anchor=west] at (-0.05,-1.95) {[tant que l'utilisateur n'est pas satisfait]};
  % messages de la boucle
  \draw[arr] (0,-2.7)  -- (4.0,-2.7)  node[midway,above,font=\scriptsize]{ajuste le brief + (Re-)Générer};
  \draw[arr] (4.0,-3.7) -- (8.0,-3.7)  node[midway,above,font=\scriptsize]{envoie le brief (prompt)};
  \draw[arr,dashed] (8.0,-4.6) -- (4.0,-4.6) node[midway,above,font=\scriptsize]{copy structuré (JSON)};
  \draw[arr] (4.0,-5.5) -- (12.0,-5.5) node[midway,above,font=\scriptsize]{copy + format};
  \draw[arr,dashed] (12.0,-6.4) -- (4.0,-6.4) node[midway,above,font=\scriptsize]{bannières (PNG)};
  \draw[arr] (4.0,-7.6) -- (0,-7.6)   node[midway,above,font=\scriptsize]{affiche l'aperçu (visualisation)};
  % --- après la boucle (satisfait) ---
  \draw[arr] (0,-9.0)  -- (4.0,-9.0)  node[midway,above,font=\scriptsize]{clique Télécharger};
  \draw[arr,dashed] (4.0,-9.9) -- (0,-9.9) node[midway,above,font=\scriptsize]{kit .zip};
\end{tikzpicture}%
}
\caption{Diagramme de séquence : génération, visualisation de l'aperçu et boucle de régénération.}
\label{fig:seq}
\end{figure}

% =====================================================================
\section{Contrat de données}
% =====================================================================
Le \emph{contrat de données} définit les structures échangées entre l'interface,
l'IA et le moteur.

\subsection{Entrée : le brief (formulaire)}
\begin{itemize}[leftmargin=1.4em]
  \item \texttt{mecanique} (1.1) : texte libre (conditions et niveaux de remise).
  \item \texttt{message} (1.2) : nom de l'opération et accroche.
  \item \texttt{illustrations} (1.3) : style graphique souhaité et éventuelle demande
        d'une image (optionnel). À défaut, la charte s'applique.
  \item \texttt{formats} : liste des bannières à produire (cases cochées).
\end{itemize}

\subsection{Sortie de Claude : le copy structuré}
Claude renvoie un \textbf{JSON} (le contrat entre l'IA et le moteur). Exemple
pour le cas de test :

\begin{lstlisting}
{
  "categorie":     "transversale",
  "titre":         "Gros projets = Grosses Remises",
  "offre_type":    "montant",
  "offre_label":   "Jusqu'a",
  "offre_nombre":  "500",
  "offre_suffixe": "EUR",
  "offre_texte":   "",
  "cta":           "J'en profite",
  "couleur_fond":  "#253081",
  "illustrer":     false
}
\end{lstlisting}

\textbf{Claude ne produit que ce JSON, jamais de HTML.} C'est le code Python qui
injecte ensuite ces valeurs dans un \textbf{gabarit fixe} (Jinja2) aux dimensions
exactes du format, avant le rendu par Playwright. Ce choix garantit l'absence
d'erreurs de parsing (HTML tronqué, guillemets mal échappés) et le respect strict
des dimensions. Les champs clés : \texttt{categorie} (une des familles produit, ou
\texttt{transversale}) sert à choisir l'image ; \texttt{offre\_type}
(\texttt{montant}, \texttt{pourcentage} ou \texttt{texte}) commande le rendu de
l'offre ; \texttt{couleur\_fond} est choisie dans la palette de la charte selon le
brief (1.3) ; \texttt{illustrer} indique si une image accompagne les pubs. Le
\textbf{même message} est conservé sur tous les formats : les gabarits adaptent la
taille du texte aux petits formats.

% =====================================================================
\section{Système de design et génération créative}
% =====================================================================
L'outil est \textbf{générique} : il ne reproduit pas un modèle unique, il produit
des bannières \textbf{variées}, adaptées à chaque brief et à chaque format. La
charte n'est pas un moule, mais un \textbf{cadre} à respecter. La conception repose
sur deux rôles complémentaires.

\begin{itemize}[leftmargin=1.4em]
  \item \textbf{Le système de design (les contraintes) :} la charte encodée en
        \emph{gabarits réutilisables} pré-construits aux dimensions exactes, avec
        \textbf{toute la palette de la charte}, la police Cabin et le logo. Chaque
        rendu est conforme à la charte. C'est le cadre non négociable qui garantit
        le respect de la marque et des dimensions.
  \item \textbf{L'IA, directrice artistique (la créativité) :} à partir du brief
        (en particulier le champ \textbf{1.3 illustrations}), Claude rédige le copy
        et \textbf{choisit la couleur de fond} dans la palette de la charte (champ
        \texttt{couleur\_fond}). Elle ne génère pas de HTML : c'est le code qui rend.
\end{itemize}

Ce choix (l'IA \emph{sélectionne et paramètre}, le code \emph{rend}) est le plus
\textbf{fiable} : aucune erreur de HTML généré, des dimensions toujours exactes, la
charte toujours respectée. La créativité vient du choix des couleurs de la charte et
de la composition, guidé par le brief.

\textbf{Style par défaut :} si le brief ne précise aucun style (champ 1.3 vide), le
fond prend le \textbf{bleu St Patricks} \texttt{\#253081}, créant une variante
purement \textbf{typographique}, très propre.

\subsection{Les types de bannières}
Chaque format appartient à un \textbf{type} qui détermine sa mise en page (le type
se déduit du nom du format) :
\begin{itemize}[leftmargin=1.4em]
  \item \textbf{Bandeau de page catégorie} (site) : image d'ambiance de la
        catégorie en fond + \textbf{carte bleue} en surimpression (titre + offre).
  \item \textbf{Header d'accueil} (site, « hero ») : grand message à gauche, offre
        (et produit détouré si le brief le demande) à droite, sur fond de la charte.
  \item \textbf{Publicité} (Google Display, Pmax, Meta) : mise en page typographique
        (accroche + offre + bouton optionnel). Si le brief (1.3) demande une image,
        le \textbf{produit détouré} est placé \emph{dans le cadre}, à côté du texte
        (disposition en ligne ou en colonne selon le format).
  \item \textbf{Formats fins} (728x90, 320x50) : gabarit dédié ; le contenu est mis
        \textbf{à l'échelle automatiquement} pour tenir, quel que soit le texte.
\end{itemize}
Les bannières \textbf{site} portent toujours une image ; les \textbf{pubs} seulement
si le brief le demande.

\subsection{La banque d'images}
Les visuels produit sont organisés par catégorie, dans
\texttt{assets/produit/categorie/<Categorie>/}, avec deux sous-dossiers :
\texttt{Bandeau page categorie} (photo d'ambiance, pour le bandeau) et
\texttt{promo} (produit \textbf{détouré}, fond transparent, pour les pubs). La
\textbf{catégorie} est déduite de l'offre par Claude ; le code va chercher la bonne
image. Pour une opération \textbf{transversale} (tout le catalogue), plusieurs
images de catégories sont assemblées \textbf{côte à côte} (montage) en fond des
bannières site.

\subsection{Une offre générique}
L'outil sert des opérations commerciales variées : l'offre peut être un
\textbf{montant} (« Jusqu'à 500 € »), un \textbf{pourcentage} (« -26 \% ») ou un
\textbf{texte} (« Livraison offerte »). Claude renseigne \texttt{offre\_type} et le
gabarit rend chacun avec le style adapté (gros chiffre, ou texte), toujours dans le
bloc jaune de la charte. Seul le \textbf{message} (l'accroche) est affiché, jamais
le détail des paliers de la mécanique, qui reste sur la page de destination.

% =====================================================================
\section{Parcours utilisateur (UX)}
% =====================================================================
La figure~\ref{fig:flow} illustre la boucle aperçu et ajustement, et la
figure~\ref{fig:wire} un schéma d'écran (wireframe).

\begin{figure}[h]
\centering
\begin{tikzpicture}[node distance=0.8cm and 1.0cm]
  \node[term] (start) {Début};
  \node[proc, right=of start] (form) {Remplir le formulaire};
  \node[proc, right=of form] (gen) {Générer};
  \node[proc, right=of gen] (prev) {Aperçu};
  \node[decision, below=1.0cm of prev] (ok) {Satisfait ?};
  \node[proc, left=1.2cm of ok] (adj) {Ajuster le brief};
  \node[term, below=1.0cm of ok] (dl) {Télécharger};

  \draw[flow] (start) -- (form);
  \draw[flow] (form) -- (gen);
  \draw[flow] (gen) -- (prev);
  \draw[flow] (prev) -- (ok);
  \draw[flow] (ok) -- node[above,font=\scriptsize]{Non} (adj);
  \draw[flow] (adj) |- (gen);
  \draw[flow] (ok) -- node[right,font=\scriptsize]{Oui} (dl);
\end{tikzpicture}
\caption{Parcours utilisateur : la boucle aperçu, ajustement, régénération.}
\label{fig:flow}
\end{figure}

\begin{figure}[h]
\centering
\begin{tikzpicture}[font=\footnotesize\sffamily]
  % cadre fenêtre
  \draw[very thick, gefblue, rounded corners] (0,0) rectangle (13,7.5);
  \fill[gefblue] (0,7.0) rectangle (13,7.5);
  \node[white, font=\small\bfseries, anchor=west] at (0.3,7.25) {Générateur de kits graphiques (Gefradis)};

  % colonne formulaire (gauche)
  \draw[gefgrey, very thick] (4.3,0.2) -- (4.3,6.8);
  \node[gefmidnight, font=\small\bfseries, anchor=west] at (0.3,6.5) {1. Le brief};
  \foreach \y/\lab in {5.7/Mécanique de remise, 4.9/Message et opération, 4.1/Style (optionnel)}{
    \node[anchor=west] at (0.3,\y+0.25) {\lab};
    \draw[fill=gefghost] (0.3,\y-0.25) rectangle (4.0,\y+0.05);
  }
  \node[gefmidnight, font=\small\bfseries, anchor=west] at (0.3,3.3) {2. Les formats};
  \foreach \y/\lab in {2.7/Site, 2.2/Google Display, 1.7/Google Pmax, 1.2/Meta}{
    \draw (0.4,\y) rectangle ++(0.25,0.25);
    \node[anchor=west] at (0.75,\y+0.12) {\lab};
  }
  \draw[fill=gefyellow] (0.3,0.35) rectangle (4.0,0.85);
  \node[black, font=\bfseries] at (2.15,0.6) {Générer le kit};

  % zone aperçu (droite)
  \node[gefmidnight, font=\small\bfseries, anchor=west] at (4.6,6.5) {Aperçu des bannières};
  \foreach \x/\y in {4.7/4.4, 7.2/4.4, 9.7/4.4, 4.7/2.0, 7.2/2.0, 9.7/2.0}{
    \draw[fill=gefblue!10] (\x,\y) rectangle ++(2.1,1.5);
  }
  \draw[fill=gefyellow] (10.0,0.35) rectangle (12.9,0.95);
  \node[black, font=\bfseries\scriptsize] at (11.45,0.65) {Télécharger .zip};
\end{tikzpicture}
\caption{Wireframe de l'interface (formulaire à gauche, aperçu à droite).}
\label{fig:wire}
\end{figure}

\subsection{Style de l'interface}
L'interface applique la charte \emph{à la manière d'un produit}, et non comme une
bannière : une \textbf{base neutre et claire} (fonds blancs et gris clairs de la
charte, \texttt{\#F9FAFB} et \texttt{\#F3F4F6}) et la marque utilisée en
\textbf{accents} (bandeau supérieur et boutons en bleu \texttt{\#253081}, bouton
d'action principal « Générer » en jaune \texttt{\#FFD52B}, titres en police Cabin).
L'objectif est un rendu d'outil professionnel et lisible, distinct du style
« bannière » réservé au livrable.

% =====================================================================
\section{Rôle de l'IA et stratégie de prompt}
% =====================================================================
\textbf{Ce que fait l'IA (Claude) :} elle joue \textbf{deux rôles}. D'abord celui
de \emph{concepteur-rédacteur} : elle rédige le copy. Ensuite celui de
\emph{directeur artistique} : elle \textbf{choisit la variante de mise en page} la
plus adaptée au brief et au format (champ \texttt{style\_variant}), et peut en
proposer plusieurs. \textbf{Ce qu'elle ne fait pas :} générer du HTML, ni produire
une image matricielle. Sa sortie est \textbf{uniquement un JSON} (copy et variante) ;
c'est le code Python qui injecte ce JSON dans un gabarit fixe, rendu en PNG par
Playwright. Cela garantit l'absence d'erreur de parsing, des dimensions toujours
exactes et le respect strict de la charte.

\textbf{Stratégie de prompt :} un \emph{prompt système} fixe le rôle et surtout les
\textbf{contraintes} (« Tu es le directeur artistique de Gefradis. Respecte
strictement la charte. Réponds uniquement par un JSON valide, sans HTML. »). Le
\emph{message utilisateur} fournit le brief et le format visé. Un exemple complet de
prompt figure en annexe (section~\ref{sec:annexe-prompt}).

\textbf{Choix : pas de génération d'image (GPT ou Imagen).} La charte (page
« Personnalité de l'imagerie ») demande des photos \emph{réelles} de produits,
brutes et épurées, et non des images inventées. Les visuels éventuels proviennent
donc des \textbf{photos réelles} de Gefradis (fond d'ambiance, produits détourés),
plutôt que d'une génération. L'architecture reste prête à accueillir un module de
génération d'illustration si le client le souhaite.

% =====================================================================
\section{Plan de construction (lotissement)}
% =====================================================================
La construction se fait par couches, de la plus simple à la plus robuste
(figure~\ref{fig:plan}) : chaque couche fonctionne avant d'attaquer la suivante,
ce qui réduit le risque sur le planning.

\begin{figure}[h]
\centering
\resizebox{\textwidth}{!}{%
\begin{tikzpicture}[font=\footnotesize\sffamily]
  \draw[arr, very thick] (-0.3,0) -- (15.3,0);
  \foreach \x/\lab in {1.2/{Étape 1\\Bannière\\statique (PNG)},
                       4.2/{Étape 2\\Tous les\\formats},
                       7.2/{Étape 3\\Intégration\\Claude},
                       10.2/{Étape 4\\Interface\\web (API)},
                       13.2/{Étape 5\\Animations\\GIF/HTML5/MP4}}{
    \fill[gefblue] (\x,0) circle (3.5pt);
    \node[proc] at (\x,1.2) {\lab};
    \draw[arr, thin] (\x,0.55) -- (\x,0.08);
  }
\end{tikzpicture}%
}
\caption{Lotissement : étapes de construction successives.}
\label{fig:plan}
\end{figure}

\textbf{Livrable cible} : une \textbf{application complète}, qui fonctionne de bout
en bout (toutes les étapes, animations comprises). C'est l'objectif, et non un
périmètre réduit. Le lotissement ci-dessus est l'\textbf{ordre de construction} :
on intègre le \textbf{flux complet} tôt (une tranche verticale sur un format),
puis on enrichit en largeur (tous les formats), en mouvement (animations) et en
finition (déploiement Cloud Run). Le module \emph{photo produit} s'ajoute selon la
réponse du client. Pour les animations, on privilégie le \textbf{GIF} (assemblé en
Python avec Pillow à partir de quelques captures PNG) ; les formats \textbf{HTML5}
et \textbf{MP4}, plus lourds, sont relégués en évolution (V2) pour ne pas mettre le
délai en péril.

% =====================================================================
\section{Décisions, hypothèses et question ouverte}
% =====================================================================
\subsection{Décisions assumées}
\begin{center}\small
\begin{tabular}{@{}>{\raggedright\arraybackslash}p{5.2cm} >{\raggedright\arraybackslash}p{9.5cm}@{}}
\toprule
\textbf{Sujet} & \textbf{Décision} \\
\midrule
Format de sortie & Images finales \textbf{non éditables} ; l'ajustement se fait en régénérant dans l'outil \\
Format statique & PNG \\
Livraison du kit & Archive \texttt{.zip} \\
Bannières animées & Plusieurs formats (HTML5, GIF, MP4) à partir d'une source unique \\
Google Pmax & Tailles standard \\
Mécanique commerciale & Donnée saisie par l'utilisateur : l'outil est \textbf{générique} \\
Référence visuelle & Fond bleu, titre blanc et chiffre clé dans un bloc jaune \\
Interface & Architecture découplée : backend \textbf{FastAPI} et frontend \textbf{HTML/Tailwind} ; UX en base neutre avec accents de marque \\
\bottomrule
\end{tabular}
\end{center}

\subsection{Question ouverte (à poser au client)}
\textbf{Les visuels (section 1.3 de la note) :} les bannières doivent-elles
intégrer des \emph{photos produit} (fenêtres, portes), ou rester \emph{basées sur
le texte, les couleurs et le logo} (sans photo) ? Si des visuels sont attendus,
une banque de \emph{photos détourées} peut-elle être fournie ?

% =====================================================================
\section{Risques et limites}
% =====================================================================
\begin{center}\small
\begin{tabular}{@{}>{\raggedright\arraybackslash}p{5.2cm} >{\raggedright\arraybackslash}p{9.5cm}@{}}
\toprule
\textbf{Risque} & \textbf{Mitigation} \\
\midrule
Complexité des bannières animées & Statique d'abord ; pour l'animé, GIF assemblé avec Pillow ; HTML5 et MP4 en V2 \\
Compatibilité Playwright et Python récent & Solution de repli (Pillow ou un rendu HTML alternatif) si besoin \\
Mémoire de Playwright (Chromium) en production & Allouer 2 à 4 Go de RAM sur Cloud Run ; limiter la concurrence des rendus \\
Assets manquants (logo source, photos) & Logo recadrable depuis la charte ; démarrage possible sans photo (fond uni) \\
Dimensions du site inconnues & À confirmer auprès du client ; valeurs standard en attendant \\
Coût et qualité du copy IA & Prompt contraint au format JSON, sorties courtes, modèle adapté au besoin \\
Complexité du frontend (FastAPI et Tailwind) & Frontend volontairement léger (pas de SPA lourde) ; construit progressivement \\
\bottomrule
\end{tabular}
\end{center}

% =====================================================================
\section{Critères de réussite et tests}
% =====================================================================
\begin{itemize}[leftmargin=1.4em]
  \item \textbf{C1 :} le rendu respecte la charte (couleurs, Cabin, logo), vérifié par contrôle visuel.
  \item \textbf{C2 :} chaque format demandé est produit aux dimensions exactes.
  \item \textbf{C3 :} le copy généré est cohérent avec le brief (cas de test Gefradis).
  \item \textbf{C4 :} la boucle aperçu, ajustement et régénération fonctionne.
  \item \textbf{C5 :} le kit se télécharge en \texttt{.zip} et s'ouvre correctement.
  \item \textbf{C6 (démonstration) :} l'outil produit un kit complet et conforme à
        la charte sur le cas de test Gefradis.
\end{itemize}

% =====================================================================
\section{Annexe : exemple de prompt}
\label{sec:annexe-prompt}
% =====================================================================
\textbf{Prompt système} (fixe, définit le rôle et les contraintes) :

\begin{lstlisting}
Tu es le directeur artistique de Gefradis (menuiserie PVC sur-mesure).
Regles STRICTES :
- Couleurs de fond autorisees : palette de la charte (#253081, #040c4d, #e7e9f4,
  #f3f4f6, #e4e4e4, #f9fafb, #fffde5) ; accent jaune #FFD52B.
- Police : Cabin uniquement.
- Garde le MEME message sur tous les formats.
- Reponds UNIQUEMENT par un JSON valide, sans aucun HTML ni texte autour.
Champs : categorie, titre, offre_type (montant/pourcentage/texte), offre_label,
offre_nombre, offre_suffixe, offre_texte, cta, couleur_fond (palette ci-dessus),
illustrer (true/false).
\end{lstlisting}

\textbf{Message utilisateur} (le brief et le format) :

\begin{lstlisting}
Operation : "Les jours Gros projets = Grosses Remises"
Mecanique : 500 EUR offerts des 1000 EUR d'achat HT
Format vise : Meta 1080x1080 (carre)
\end{lstlisting}

\textbf{Sortie attendue} (JSON seul, sans HTML) :

\begin{lstlisting}
{
  "categorie": "transversale",
  "titre": "Gros projets = Grosses Remises",
  "offre_type": "montant",
  "offre_label": "Jusqu'a",
  "offre_nombre": "500",
  "offre_suffixe": "EUR",
  "offre_texte": "",
  "cta": "J'en profite",
  "couleur_fond": "#253081",
  "illustrer": false
}
\end{lstlisting}

\vspace{1em}
\hrule
\vspace{0.5em}
{\small\itshape Fin du cahier de conception (version 0.3). Ce document est destiné
à évoluer selon la réponse du client à la question ouverte et selon les retours
sur les premières versions. Version 0.3 : ajout des types de bannières, de la
banque d'images, du montage transversal et de l'offre générique.}

\end{document}
